% ============================================================
% frontmatter/titlepage_cn.tex —— 中文题名页（规范 3.4）
% 【V4 校准】严格按“附件1-11 p.1（学术型硕士式样）”重建：
%   顺序（顶→底）：
%     学校代码/研究生学号 → 中国地质大学 → 硕士学位论文
%     → 论文题目 → 姓名 → 学科专业 → 指导教师 → 培养单位 → 完成日期
%   此前（V3）误用了师兄“硕士专业学位”题名页的顺序，与学术型附件相反，
%   本次已纠正为附件学术型顺序。
%   字体：校名/类别 宋体一号加粗；论文题目 黑体二号加粗；信息块 宋体三号；
%        学号等阿拉伯数字用 Times New Roman（\tnr）。
%   标签经 \labelfour 等宽分散（与附件式样一致：2 字大间距、4 字小间距）。
%   注：沿用 tabular（fixLineskip 修复后正常分行、顺序正确），不塌行。
% ============================================================
\clearpage
\thispagestyle{empty}

\begin{center}
  % —— 整体下移，使顶部“学校代码”落在距页顶约 5.8cm（对齐附件学术型式样）——
  \vspace*{2.8cm}
  % —— 顶部：学校代码 / 研究生学号（左锚 r0≈0.019）——
  \zihao{3}\songti
  \makebox[\textwidth][l]{%
    \hspace*{0.019\textwidth}%
    学校代码：\tnr{\thesisSchoolCode}\quad 研究生学号：\tnr{\thesisStudentID}}

  \vspace*{1.1cm}
  % —— 学位授予单位名称（宋体一号加粗居中）——
  {\songti\zihao{1}\textbf{\thesisSchool}}

  \vspace*{1.1cm}
  % —— 学位论文类别（宋体一号加粗居中）——
  {\songti\zihao{1}\textbf{\thesisCategory}}

  \vspace*{1.3cm}
  % —— 论文题目（黑体二号加粗居中，超宽自动折行；块宽约 0.70 版心使两行均衡）——
  \parbox{0.70\textwidth}{\centering\heiti\zihao{2}\thesisTitle}

  \vspace*{1.7cm}
  % —— 信息块（左锚 r0≈0.268；顺序按附件：姓名→学科专业→指导教师→培养单位）——
  \makebox[\textwidth][l]{%
    \hspace*{0.268\textwidth}%
    \begin{tabular}{rl}
      \labelfour{姓名}：     & \zhlabelsep\thesisAuthor   \\[0.42cm]
      \labelfour{学科专业}： & \zhlabelsep\thesisMajor   \\[0.42cm]
      \labelfour{指导教师}： & \zhlabelsep\thesisAdvisor \\[0.42cm]
      \labelfour{培养单位}： & \zhlabelsep\thesisOrg
    \end{tabular}}

  \vspace*{3.2cm}
  % —— 完成日期（汉字，居中）——
  \thesisDate
\end{center}
